\clearpage
\renewcommand{\footerfilename}{SystemAdministration.tex}

\section{System Administration} \label{sec:SystemAdministration}

\subsection{Set up the machine}

For Ubuntu, the packages ``iraf'' and ``python3-pyraf'' need to be installed and
ready to go. It may be better to install the PyRAF from the github repository.

This basic setup will be augmented.

It is necessary to create a {\textasciitilde{}/.iraf} directory and
add the git clone's loginuser.cl and your custom
pyraflogin3.py\footnote{\url{https://github.com/The-SMTSci/SMTSCI-iraf.git}}
programs. These are biased towards reducing longslit small-telescope
spectra.

Create a {\textasciitilde{}/bin} directory and populate with the bin
files from the git clone.

The {\textasciitilde{}/bin} directory has to be included with the path.

When Unix starts a bash shell, it moves to the home directory you have
registered when the account was created (/etc/passwd file). It looks
for initialization files, like \dhl{.bashrc}. In that file, it tests to
see if you created a \dhl{.bash.aliases} file -- and if found it will
'source' that file. Sourcing the file means all the action, like
exporting variables and declaring functions will in inserted into
the current shell. In your .bash.aliases file, you can add your
own tests for additional files (like \dhl{.pyraf.aliases}!) and it
will source them as well.

From the SMTSCI-iraf.git environment, copy the dot.bash\_aliases to
{\textasciitilde{}/.bash\_aliases} and copy the dot.pyraf.aliases to
{\textasciitilde{}/.pyraf.aliases}' (Git does not allow for hidden
files to be saved).

Reduction steps discussed here will use at least these scripts:

\begingroup \fontsize{10pt}{10pt}
\selectfont
%%\begin{Verbatim} [commandchars=\\\{\}]
\begin{verbatim} 
fixnames
fitserial
trim
\end{verbatim}
\endgroup
%% \end{Verbatim}
fitsls
